\section{The explicitness check}\label{sec:xpl}

This appendix describes the Lean command \texttt{\#check\_explicit}, which checks the syntactic form of level L1
(\Mref{Definition}{def:intro-levels}) on the actual construction terms. It complements the kernel's axiom report: the
kernel checks that a theorem proves its stated proposition, while the command audits the data dependencies of a
definition. It is not used in any proof of this supplement. Its implementation is in
\texttt{SymplecticPacking/Explicit.lean}.

\subsection{Data dependencies}\label{sec:xpl-traversal}

Starting at the named definition, the command follows constants appearing in
data positions, including the bodies of lambda expressions and local values.
It skips theorem bodies and arguments whose types are propositions, so a
classical existence proof that justifies a formula does not make the formula
an unspecified choice; a choice that supplies a value is inspected.

Project and vendored definitions are traversed. The audit rejects direct witness
selection by \texttt{Classical.choose}, \texttt{Function.invFun} and the other
recorded choice operators unless it belongs to an accepted case below.
A missing proof or value represented by \texttt{sorryAx} is rejected.
The source scan separately excludes new axioms and unchecked implementation
mechanisms. The kernel audit of the main results permits only
\texttt{propext}, \texttt{Classical.choice} and \texttt{Quot.sound}.

\subsection{Definite operations and searches}\label{sec:xpl-accept}

\texttt{Explicit.the} supplies definite descriptions with proofs of uniqueness.
Its uses in the construction specify integral
curves, inverse images, scalar roots and linear solves. Each defining predicate
and its domain are part of the mathematical construction. Classical decisions
on propositions are allowed as case distinctions at level L1; they do not
establish computability from real names.

For a search over naturals, rationals or finite encoded data, the command checks
both the search type and genuine decidability of its predicate. A decidability
argument reaching classical propositional decision is rejected. Requirements
on a wrapper's arguments are propagated to its call sites, so wrapping a search
in another definition does not remove the requirement. Search predicates on
function spaces are rejected. The implementation also checks natural infima,
filtered finite sets and definite descriptions over search types.

Mathlib uses classical choice to define limits, integrals and other definite
mathematical operations. The audit accepts only the exact reviewed declarations
listed in \texttt{Explicit/}\allowbreak\texttt{Whitelist.lean}; it records their dependency paths
in its output. These entries are library foundations. They do not permit witness
selection in a project or vendored construction. The finite charts, primitives,
cutoffs and flows of the construction are therefore audited through their actual
data dependencies.

\begingroup\sloppy
\subsection{Definite descriptions in the construction}\label{sec:xpl-descriptions}

The descriptions over continuous types that the command accepts are reviewed by hand (\Msec{sec:ver}). The report of
\texttt{\#check\_explicit} on \texttt{General.packingHamiltonian} lists $26$ declarations containing
\texttt{Explicit.the}. Each describes one of the following objects; in every case the defining property is an
explicit equation or inequality in the data of the construction, uniqueness is proved, and none refers to a packing
specification.
\begin{itemize}
\item \emph{Inverses of injective maps} (the unique preimage, with a fixed fallback value off the image):
 \texttt{Inner.invR}, \texttt{Inner.TrotInv}, \texttt{Qhat.ordInv}, \texttt{Qhat.qhatInv},
 \texttt{HighDim.ToricPackingStages.circleClockInv}, \texttt{layerChartInv} and \texttt{nativeInverse},
 \texttt{BranchLarge3.CubicBridge.expInv} (inverse of a radial moment map), and \texttt{General.Ham.pushforward}
 (the unique source point $u\in\EB(r)$ with $F_t(u)=y$, Definition~\Mnum{def:ham-H}).
\item \emph{Unique roots and critical points}: \texttt{Strip.layerRoot} and \texttt{Strip.Cubic.layerRoot} (the layer roots
 of Lemma~\ref{lem:str-layer-root} and of the four-dimensional layer maps of \S\ref{sec:shd-four}), \texttt{Strip.vOfT}, \texttt{BranchLarge3.PrimExt.rootExt} (the
 root closest to an explicit centre), \texttt{Branch2.Planar.PlanarInput.pstar} and
 \texttt{HighDim.PlanarGen.Input.vstar} (the unique critical point of a strictly convex potential).
\item \emph{Solutions of linear equations}: \texttt{BranchLarge3.TargetMoser.solveL} (an invertible bilinear form),
 \texttt{BranchSmall.VT.transportField} and \texttt{returnField} (the Moser fields defined by nondegenerate linear
 systems, \S\ref{sec:vtr}), and \texttt{General.Ham.Linear.posSqrt} (the positive square root of a positive
 semidefinite matrix).
\item \emph{Solutions of ordinary differential equations} (the unique integral curve of a Lipschitz field):
 \texttt{Qhat.flowOf}, \texttt{Return.flowFrom}, \texttt{BranchSmall.VT.flowI}, \texttt{BranchLarge3.UMoser.curve},
 \texttt{HighDim.ToricPackingStages.nativeAmbientFlow}.
\item \emph{Extrema and measures}: \texttt{HighDim.ToricPackingStages.finiteMinimumOrZero} (the least value of finitely
 many continuous functions on compact sets, $0$ if there is none) and \texttt{Explicit.planeVolume} (Lebesgue measure
 on $\R^2$, characterised by its values on rectangles).
\end{itemize}

\subsection{Vendored declarations used}\label{sec:xpl-vendored}

The $3810$ declarations of the vendored library in the dependency closure of the statement and algorithm theorems
(\Msec{sec:ver}) are listed by module in \texttt{verification/Audit/VendoredUsage.lean}. They are explicit charts and
coordinates on the diagonal quadrics and the Fermat cubic, Fubini--Study and Liouville primitives, toric and Hopf
coordinates, Moser fields and their naturality under pullback, planar area and rearrangement bounds, and the moment
polynomials of the models. The closure contains nothing from the curve-analysis directories \texttt{Compactness}, \texttt{Continuation},
\texttt{Degree}, \texttt{Cauchy} and \texttt{Campanato}.

\endgroup

\subsection{Formal statements that differ in form from the text}\label{sec:xpl-interface}

Table~\ref{tab:xpl-interface} records the statements of the paper and of this supplement whose Lean counterparts have a
different form. The complete printed statements are proved in the text; where Lean records only the consequence
needed by the construction, the statement carries the visible label \emph{needed consequence verified}.
The label \emph{full statement verified} is reserved for full mathematical counterparts, up to the stated
presentation differences. Unqualified source tags give references without a full-coverage claim. The scope labels
record a human comparison, not a kernel check of the prose. Lean names are relative to \texttt{SymplecticPacking}.

{\footnotesize
\renewcommand{\arraystretch}{1.12}
\begin{longtable}{@{}p{0.22\textwidth}p{0.74\textwidth}@{}}
\caption{Formal statements that differ in form from the text.}\label{tab:xpl-interface}\\
\toprule
statement & formal form\\\midrule
\endfirsthead
\multicolumn{2}{c}{\tablename\ \thetable\ (continued)}\\
\toprule
statement & formal form\\\midrule
\endhead
\midrule
\multicolumn{2}{r}{Continued on the next page}\\
\endfoot
\bottomrule
\endlastfoot
flat step $S$ (Def.~\ref{def:pre-flat}) & Mathlib's \texttt{Real.smoothTransition}, built from \texttt{expNegInvGlue}${}=a_\flat$; the
 definite description $\defdesc(h)$ is \texttt{Explicit.the}\\
real coordinates & the coordinate $\zeta_l$ of $\R^{2n}$ (Table~\ref{tab:pre-interfaces}) is the coordinate of index $l-1$ of
 \texttt{Problem.E}~$n$; the coordinates are called $(q,p)$ there, with $p$ a coordinate and not a moment\\
Thm.~\Mnum{thm:cmp-realisation} (Thm.~\ref{thm:plc-lift}) & the formal statement records (c) with real data, that every
 $q^{(j)}_t$ maps $V_\Delta$ into itself, and, in place of (b), one collar width on which $Q^{(j)}_t$ Poisson-commutes with
 every face moment; the remaining assertions of (a) and (b) are proved in Theorem~\ref{thm:plc-lift} only\\
Thm.~\ref{thm:pln-isotopy} & stated for a single ordinary path, without the parameter $y$, as used in dimension four; the
 smooth dependence of the planar flow on a finite-dimensional parameter is a separate statement,
 \texttt{HighDim.PlanarGen.compact\_param\_flow\_smooth}\\
Cor.~\ref{thm:pln-full-disc} & the case of Theorem~\ref{thm:pll-inner}(e): $g(I)=I+4e$, $\eps=0$, $\Upsilon\le h$ on $[0,h]$, and the
 core hypothesis $\Upsilon(\pi|\mathfrak e_1(z)|^2)<\pi|z|^2+\frac{13}{24}e$ on $\bar D(h-4e)$\\
Def.~\ref{def:shd4-data} & valid cubic data admit for $(q_t)$ only the whole-disc isotopy $\hat q$ of an ordinary path of disc
 area $h$, which both four-dimensional branches use\\
Prop.~\ref{prop:shd4-moser} & items (a)--(d) and the smoothness of $\Phi$ and $\Phi^{-1}$; the decomposition
 $\Phi=\Phi_2\circ\Phi_1\circ\Phi_0$ is not formalised\\
Prop.~\ref{prop:shd4-maps}(c) & in the form used in Theorem~\ref{thm:cub-packing}: if $\mathcal Tf^N_i(X)=\mathcal Tf^N_j(X')$ for the
 transfer $\mathcal T$ of Theorem~\ref{thm:rtc-transfer}, then $f^N_i(X)=f^N_j(X')$\\
Prop.~\ref{prop:rtc-return}, Lem.~\ref{lem:rtc-ball-chart}, Thm.~\ref{thm:rtc-transfer} & there is no single statement about $\mathcal T$ on its whole domain; the
 theorem is formalised along each embedding $\mathcal T\circ f^N_i$. The pullback identities are stated along smooth families
 of unit vectors (return and ball chart) and for $\mathcal T_{\rm u}$ composed with smooth families with values in $\mathcal N_h$,
 together with Proposition~\ref{prop:rtc-central}; injectivity is stated as in the previous row\\
Thm.~\ref{thm:cub-packing} & for the physical maps $u\mapsto\sqrt C\,f_i(u/\sqrt C)$: smoothness with symplectic derivative on the
 open domain $\Bo(C\rho_i^{\rm dom})$, injectivity on $B(C\rho_i^+)$, pairwise disjoint images of these closed balls, and images in
 $\Bo(C)$; the margin $m$ is not stated, and (iv) is structural. Injectivity, disjointness and $\pi|f_i|^2<1$ on all of
 $\Bo(\hat\rho_i)$ are proved for the capacity-scale maps in both cases\\
Rem.~\ref{rem:rtc-equal} & the equal-capacity data form the structure \texttt{Params}, whose field \texttt{Params.}$\varepsilon$ is $4\eta_\star$ (not the
 $\eps$ of \Mref{Theorem}{thm:intro-4d}) and whose \texttt{Params.rho}${}=(1-\eta_\star)/\sqrt k$ is $\rho$; \texttt{Branch2.CubicData.ofParams} turns them into
 cubic data, so both cases pass through \texttt{Branch2.cubicEmbeddings\_spec}, the map $u\mapsto\sqrt C\,f_i(u/\sqrt C)$. The data of
 \Mref{Theorem}{thm:intro-4d} are those of \texttt{General.epochHamiltonianQ}, and \texttt{General.epochQ\_equalCase} proves
 \texttt{Branch2.EqualCase} for them\\
support of $H$ (\Msec{sec:ham}) & \texttt{General.Ham.isotopySupport} is the union of
 $F_{i,t}(B(r_{2,i}))$ over $i$ and $t\in[0,1]$; \texttt{isotopySupport\_compact} proves compactness and
 \texttt{ofIsotopy\_support\_precise} proves that $H$ vanishes outside it, under the common-isotopy hypotheses.
 For $H[n,k,C,c,s]$, \texttt{General.packingSupportSet\_compact},
 \texttt{packingHamiltonian\_support\_precise} and \texttt{packingSupportSet\_subset\_radius} give the compact
 swept set $K_H$ and its containment in the closed ball of radius \texttt{General.packingSupportRadius}.
 The resulting vanishing outside that ball is also stated by \texttt{main\_support\_radius}\\
\end{longtable}
}

\subsection{Reproduction}\label{sec:xpl-searches}

The statement file runs the check on the Hamiltonian of \Mref{Theorem}{thm:main} and on the Hamiltonian of the Epoch
problem (\Mref{Remark}{rem:eff-epoch-interface}); further checks inspect each embedding constructor and, for the finite
output, the polynomial interpreter, the rational verifier and the output construction. A nonempty violation list fails
the build. Besides the rejected searches listed in \Msec{sec:ver}, the regression file \texttt{test/ExplicitTest.lean}
contains genuine finite rational searches that must be accepted, a natural infimum, filtering and a definite
description. These tests exercise the search rule; the soundness and termination of the finite packing search are
proved in \S\S\ref{sec:certificate} and~\ref{sec:alg}.
